\customchapter{CONCLUSIONES}

Las conclusiones se enuncian en correspondencia con los dos objetivos específicos
que incorporan trabajo computacional, precedidas del diagnóstico que fundamenta
el proyecto y seguidas de la lectura analítica que los integra.

\begin{enumerate}
\item \textbf{Diagnóstico.} La clasificación de elementos estructurales sobre
levantamientos TLS de patrimonio arquitectónico carece hoy de solución
automatizada. La segmentación semántica 3D de conjunto cerrado no puede cubrirla
porque fuerza todo elemento no previsto hacia una clase conocida, y el marco que
sí levanta ese supuesto ---NCD--- rinde entre 11,8 y 47,5 mIoU en clases
novedosas, por debajo del umbral de utilidad práctica. La revisión crítica
localiza la causa plausible en el descriptor geométrico calculado a mano del
módulo de prototipos, y la reproducción de la línea base
(Sección~\ref{sec:linea-base}) fija el punto de partida sobre datos reales, con
el resultado negativo del cambio de normalización descartando una explicación
alternativa.

\item \textbf{Conclusión asociada a OE1 (integración).} La sustitución de
$f_{\mathrm{geo}}$ por la representación auto-supervisada congelada
$f_{\mathrm{ssl}}$ es \emph{[viable / viable bajo condiciones / no viable]}
dentro del módulo de prototipos de HyperNCD, conservando intacto el razonamiento
por hipergrafo. La compatibilidad de resolución entre el \emph{encoder} PTv3 y la
rejilla de Minkowski se resuelve mediante \emph{up-cast} y \emph{cache} por
sección, y la prueba de funcionamiento sobre la sección de control confirma que
el mIoU \emph{known} no se degrada respecto de la línea base.

\item \textbf{Conclusión asociada a OE2 (calibración del híbrido).} El parámetro
que domina el comportamiento del modelo combinado es \emph{[por determinar]}; el
valor de $\alpha$ que maximiza el mIoU \emph{novel} indica cuánto peso conserva
la componente geométrica una vez reemplazado el descriptor manual. La atribución
por etapas muestra que la ganancia se origina principalmente en \emph{[la
representación / el razonamiento por hipergrafo]}, resultado que no sería
observable si el híbrido se evaluara como caja negra.

\item \textbf{Lectura integrada (OE3).} La mejora se cuantifica mediante
$\Delta$mIoU \emph{novel} y la brecha cerrada $G$, reportando siempre los tres
términos que componen esa métrica y la clase que domina el agregado. El diseño
aísla el reemplazo en un único punto de fusión, lo que hace que cualquier
variación de la métrica sea atribuible al cambio y convierte el estudio de
ablación en la evidencia central del trabajo.

\item \textbf{Alcance y límites.} El estudio se realiza sobre un edificio
nombrado y un levantamiento de $\sim$1--2\,GB organizado en cinco secciones, con
rangos de operación declarados por adelantado y verificados empíricamente. El
proyecto está limitado por cómputo y su ejecución depende del hardware
efectivamente asignado en el clúster Khipu, cuya configuración se verifica y no
se estima, lo que hace el proyecto ejecutable con los recursos disponibles y
verificable en sus límites.
\end{enumerate}
